\documentclass[../main]{subfiles}
\begin{document}
\newpage

\subsubsection{\texttt{bj} --- Box--Jenkins}

\begin{lstlisting}
theta, m = pysib.bj(u, y, nb, nc, nd, nf, nz)
\end{lstlisting}

Identifies a Box--Jenkins model by minimizing the prediction-error
criterion.  The model structure is

\[
  y(t) = \frac{B(q^{-1})}{F(q^{-1})}\,u(t - n_k)
       + \frac{C(q^{-1})}{D(q^{-1})}\,e(t),
\]

where $e(t)$ is white noise.
The one-step-ahead prediction error is

\[
  \varepsilon(t;\theta) =
  \frac{D(q^{-1})}{C(q^{-1})}
  \left[y(t) - \frac{B(q^{-1})}{F(q^{-1})}\,u(t-n_k)\right],
\]

and the criterion $V(\theta) = \tfrac{1}{N}\sum_t \varepsilon(t;\theta)^2$
is minimized.  The problem is nonlinear in $\theta$ because $C$, $D$,
and $F$ all appear in denominators.
An ARX estimate provides the initial guess for $B$ and $F$; $C$ and
$D$ are initialized from the ARX denominator.  The optimizer uses
steepest-descent and Newton steps implemented in a C extension, with
LAPACK used for the Newton step.

\subsubsection*{Arguments}
\begin{center}
\begin{tabular}{lll}
  \toprule
  Name & Type & Description \\
  \midrule
  \texttt{u}  & array\_like & Input signal. \\
  \texttt{y}  & array\_like & Output signal. \\
  \texttt{nb} & int & Number of $B$ parameters. \\
  \texttt{nc} & int & Number of $C$ parameters. \\
  \texttt{nd} & int & Number of $D$ parameters. \\
  \texttt{nf} & int & Number of $F$ parameters. \\
  \texttt{nz} & int & Input delay $n_k \geq 0$. \\
  \bottomrule
\end{tabular}
\end{center}

\subsubsection*{Returns}
\begin{center}
\begin{tabular}{lll}
  \toprule
  Name & Type & Value \\
  \midrule
  \texttt{theta} & ndarray & $[b_0,\ldots,b_{n_b-1},\; c_1,\ldots,c_{n_c},\;
                              d_1,\ldots,d_{n_d},\; f_1,\ldots,f_{n_f}]$ \\
  \texttt{m}     & dict    & \texttt{B}, \texttt{C}, \texttt{D}, \texttt{F} free;
                             \texttt{A}$=[1]$ \\
  \bottomrule
\end{tabular}
\end{center}

\end{document}
